\documentclass[11pt,a4paper]{article}

% --------------------
% Pacchetti
% --------------------
\usepackage[italian]{babel}
\usepackage[T1]{fontenc}
\usepackage[utf8]{inputenc}
\usepackage{geometry}
\geometry{margin=2.5cm}

\usepackage{amsmath, amssymb}
\usepackage{booktabs}
\usepackage{array}
\usepackage{tabularx}
\usepackage{enumitem}
\usepackage{graphicx}
\usepackage{multicol}
\usepackage{setspace}

\setstretch{1.15}

% --------------------
% Comandi utili
% --------------------
\newcommand{\checkbox}{\(\square\)}

% --------------------
% Documento
% --------------------
\begin{document}

\begin{center}
\textbf{\Large Liceo Scientifico ``Galileo Galilei'' -- San Donà di Piave (VE)}\\[0.2cm]
\textbf{\large Laboratorio di Fisica -- Scheda di Progettazione}\\[0.4cm]
\end{center}

\noindent
\textbf{Classe:} II / III \hfill
\textbf{Durata:} 1 ore di laboratorio + 1 ora di rielaborazione dati + redazione della relazione.

\vspace{0.4cm}

\section*{Titolo}
\textbf{Analisi del moto parabolico attraverso la video--analisi con Tracker}

\section*{Obiettivi di apprendimento}
Al termine dell’attività lo studente sarà in grado di:
\begin{itemize}
    \item riconoscere le caratteristiche del moto parabolico come combinazione di due moti indipendenti:
    \begin{itemize}
        \item moto rettilineo uniforme (componente orizzontale);
        \item moto uniformemente accelerato (componente verticale);
    \end{itemize}
    \item utilizzare il software \textit{Tracker} per analizzare un video reale e ricavare dati cinematici
    (posizione, velocità, accelerazione);
    \item confrontare i dati sperimentali con i modelli teorici;
    \item sviluppare capacità di lavoro cooperativo e di comunicazione scientifica sintetica;
    \item riflettere sul proprio percorso di apprendimento (autovalutazione).
\end{itemize}

\section*{Strumenti e materiali}
\begin{itemize}
    \item PC o notebook con \textit{Tracker} installato;
    \item video preregistrato o registrato in classe di un moto parabolico (es. pallina lanciata in orizzontale);
    \item fogli di calcolo (Excel, Google Sheets) per l’elaborazione dei dati;
    \item scheda di progetto e scheda di autovalutazione.
\end{itemize}

\section*{Fasi del lavoro}

\subsection*{1. Creazione del video sperimentale (20 minuti)}
L’attività laboratoriale inizia con la produzione autonoma del video del moto parabolico, che costituirà il dato sperimentale da analizzare.

\begin{itemize}
    \item Registrazione del moto tramite \textbf{telefono cellulare}.
    \item Trasferimento del file video sul PC per l’analisi tramite il software \textit{Tracker}.
\end{itemize}

Durante la fase di registrazione del video è necessario rispettare le seguenti accortezze sperimentali:

\begin{itemize}
    \item \textbf{Definizione dell’unità di misura:} fissare sul muro (preferibilmente bianco) una scala metrica utilizzando scotch di carta, in modo da disporre di un riferimento per la calibrazione delle distanze.
    \item \textbf{Condizioni del moto:} lanciare la pallina in modo da realizzare un moto parabolico ben visibile, evitando urti o rotazioni anomale.
    \item \textbf{Qualità della registrazione:} mantenere il cellulare ben fermo durante la ripresa, evitando movimenti o vibrazioni che possano compromettere l’analisi.
    \item \textbf{Inquadratura:} assicurarsi che l’intera traiettoria del moto rientri nel campo visivo della videocamera.
\end{itemize}

\subsection*{2. Analisi del moto tramite Tracker (30 minuti)}
Una volta acquisito il video, si procede all’analisi quantitativa del moto mediante il software \textit{Tracker}.

\begin{itemize}
    \item Importazione del video all’interno di \textit{Tracker}.
    \item Definizione del sistema di riferimento cartesiano.
    \item Calibrazione delle distanze utilizzando l’unità di misura fissata sul muro.
    \item Marcatura della traiettoria punto per punto.
    \item Raccolta dei dati cinematici: coordinate $x$, $y$ e tempo.
\end{itemize}

Analisi dei grafici forniti dal software:
\begin{itemize}
    \item $x(t)$: verifica della linearità della componente orizzontale (moto rettilineo uniforme);
    \item $y(t)$: verifica della legge oraria del moto uniformemente accelerato;
    \item traiettoria $y(x)$: verifica dell’andamento parabolico.
\end{itemize}

\subsection*{3. Analisi dei dati e rielaborazione di gruppo (2 ore in classe)}
I dati ottenuti tramite il software \textit{Tracker} vengono esportati e rielaborati dal gruppo di lavoro al fine di confrontare il modello teorico con il comportamento sperimentale del sistema.

\begin{itemize}
    \item Esportazione dei dati cinematici ($x$, $y$, $t$) da \textit{Tracker} in un foglio di calcolo
    (Excel, Google Sheets o software equivalente).
    \item Rielaborazione dei dati:
    \begin{itemize}
        \item organizzazione delle misure in tabelle;
        \item costruzione dei grafici $x(t)$, $y(t)$ e $y(x)$;
        \item individuazione dei parametri caratteristici del moto.
    \end{itemize}
    \item Selezione di almeno \textbf{quattro punti} significativi della traiettoria.
    \item Calcolo dei valori teorici corrispondenti sulla base delle leggi del moto parabolico.
    \item Confronto tra valori sperimentali e teorici.
    \item Calcolo e discussione degli scostamenti percentuali, con particolare attenzione alle possibili fonti di errore sperimentale.
\end{itemize}


\subsection*{4. Redazione della relazione di laboratorio (2 ore)}
Al termine dell’analisi dei dati, ogni gruppo procede alla redazione di una relazione di laboratorio strutturata, seguendo lo schema condiviso dal docente.

La relazione dovrà documentare in modo chiaro e rigoroso l’intero percorso svolto, dal contesto sperimentale all’interpretazione dei risultati, evidenziando il confronto tra modello teorico e dati sperimentali.

In particolare, la relazione dovrà includere:
\begin{itemize}
    \item \textbf{Introduzione e obiettivo dell’esperienza:} descrizione del fenomeno studiato e finalità dell’attività.
    \item \textbf{Descrizione dell’apparato sperimentale e della procedura:} modalità di acquisizione del video e condizioni sperimentali adottate.
    \item \textbf{Analisi dei dati:} presentazione delle tabelle e dei grafici ottenuti tramite il foglio di calcolo.
    \item \textbf{Confronto teorico--sperimentale:} validazione del modello del moto parabolico attraverso il confronto tra valori teorici e sperimentali.
    \item \textbf{Discussione degli errori:} analisi critica degli scostamenti e delle principali fonti di incertezza.
    \item \textbf{Conclusioni:} sintesi dei risultati e riflessione sul significato fisico dell’esperienza.
\end{itemize}



\end{document}
